% TOP_PROBLEM: {"problemId": "problem.coding-theory-and-combinatorics-problem-cc-6-existence-of-finite-projective-planes-of-non-prime-power-orders", "canonicalTitle": "Coding Theory and Combinatorics Problem CC-6: Existence of Finite Projective Planes of Non-Prime-Power Orders", "releaseRank": 404, "primaryDomain": "cryptography_coding_information_optimization", "primaryDomainLabel": "Cryptography, coding, information & optimization", "displayStatus": "open", "releaseStatus": "open"}
% !TeX program = lualatex
% Definition and sources only; this file is not an LLM solution attempt.
\documentclass[11pt]{article}
\usepackage[margin=25mm]{geometry}
\usepackage{fontspec}
\setmainfont{Latin Modern Roman}
\usepackage{amsmath,amssymb,mathtools}
\usepackage{unicode-math}
\setmathfont{Latin Modern Math}
\usepackage{xurl,hyperref}
\hypersetup{colorlinks=true,urlcolor=blue}
\setcounter{secnumdepth}{0}
\setlength{\parindent}{0pt}
\setlength{\parskip}{5pt}
\setlength{\emergencystretch}{3em}
\begin{document}
\section*{\texorpdfstring{Coding Theory and Combinatorics Problem CC-6: Existence of Finite Projective Planes of Non-Prime-Power Orders}{Coding Theory and Combinatorics Problem CC-6: Existence of Finite Projective Planes of Non-Prime-Power Orders}}\label{404}
\subsection{Definitions and mathematical statement}
A finite projective plane is a finite point--line incidence structure in which two distinct points lie on a unique line, two distinct lines meet at a unique point, and four points exist with no three collinear. Its order $q$ means each line contains $q+1$ points; it then has $q^2+q+1$ points and as many lines. The problem asks whether a plane exists of order $12$, or more generally of any order $q\ge2$ which is not a prime power and is not excluded by the Bruck--Ryser--Chowla obstruction. In particular, for $q\equiv1,2\pmod4$ that obstruction requires $q$ to be a sum of two integer squares. Satisfying the obstruction is only a necessary condition, not a construction. The order-12 question is a particular case of the broader existence question.
\par\smallskip\textit{Formulation and concept references:} \href{https://arxiv.org/pdf/0811.1254}{[S1]}.

\subsection{Short English statement}
Can a finite projective plane have order twelve, or any other allowable order that is not a prime power?
\subsection{Sources}
\begin{itemize}
\item[S1]Huber, Coding theory and algebraic combinatorics (2008), Remark 1.1.
\url{https://arxiv.org/pdf/0811.1254}.
\textit{Formulation source.}
\end{itemize}
\end{document}
